In this section, we discuss agents with inaccurate belief. \Cref{thm:inaccurate_belief} gives bounds on the utility lost as a result of the agent having an inaccurate belief. We give an upper bound in terms of the (weaker) absolute error and lower bounds in terms of both absolute and relative error, showing that the upper bound is tight up to factors of $2$ and $4$ for absolute and relative error.

\begin{theorem} \label{thm:inaccurate_belief}
Let $A$ be a perfect expected utility maximizer whose belief $\rho$ has absolute error $\epsilon$ with respect to the correct belief $\rho^*$. Then it is a $(\frac{2}{1-\gamma} - \frac{2}{1- \gamma(1-\epsilon)})$-optimizer with respect to $\rho^*$ and this bound is tight up to a factor of $2$. Moreover, if $\epsilon$ is the relative error, this bound is tight up to a factor of 4.
\end{theorem}

So far, we have considered the worst-case scenario. In the next theorem, we show that in the case of both absolute and relative error, the upper bound is tight up to a constant factor even in the case when the error at each timestep is randomly chosen from the set $\{-\epsilon, \epsilon\}$ (that is, $\rho(e_t|\mae_{<t}a_t)$ is chosen uniformly from the set $\{\max(0,\rho^*(e_t|\mae_{<t}a_t)-\epsilon), \min(1,\rho^*(e_t|\mae_{<t}a_t)+\epsilon)\}$ in the case of absolute error and $\{\frac{\rho^*(e_t|\mae_{<t}a_t)}{1+\epsilon}, \min(1,(1+\epsilon)\rho^*(e_t|\mae_{<t}a_t))\}$ in the case of relative error), independently of other timesteps.

\begin{theorem} \label{thm:avg_case_belief}
Let $A$ be a perfect expected utility maximizer whose belief $\rho$ is such that for any $e_t,\mae_{<t}a_t$, the value $\rho(e_t|\mae_{<t}a_t)$ is independently for any argument chosen uniformly from the set $\{\max(0,\rho^*(e_t|\mae_{<t}a_t)-\epsilon), \min(1,\rho^*(e_t|\mae_{<t}a_t)+\epsilon)\}$ in the case of absolute error and $\{\frac{\rho^*(e_t|\mae_{<t}a_t)}{1+\epsilon}, \min(1,(1+\epsilon)\rho^*(e_t|\mae_{<t}a_t))\}$ in the case of relative error.

Then, in expectation, the amount of expected discounted utility lost is respectively
\begin{gather}
\frac{1}{1-\gamma} - \frac{1}{1- \gamma(1-\epsilon/8)}\\
\frac{1}{1-\gamma} - \frac{1}{1- \gamma(1-\epsilon/16)}
\end{gather}

and this is equal to the upper bound for non-random error up to a factor of $16$ and $32$, respectively.
\end{theorem}

